% ======================================================================
% SECTION 6: CONCLUSIONS
% Four thematic blocks expanded from the v0.8.0 bracketed draft at
% 2030-pdac-1min/paper/draft-paper/sections/conclusions.tex.
% Location: 2030-pdac-1min/paper/full-paper/sections/conclusions.tex
% ======================================================================

\subsection{Headline Artifact Result}
\label{sec:conclusions-headline}

The headline artifact of this v0.9.0 release is a complete five
phase Claude Code Opus 4.7 1M Max workflow (code instructions plus
code generation plus execution plus a draft template, now expanded into a full paper, followed by this final version) that produces the first 60 second 8 arm Whipple procedure plus
Daraxonrasib precision oncology adjuvant simulation in the
literature.\ The headline numbers tie back to file paths located in the repository under
\texttt{2030-pdac-1min/paper/execution/}: mean composite of 93.298
across 32 iterations with std 1.225 and 95\% CI half width of 0.462 at
\texttt{2030-pdac-1min/paper/execution/iterations/iteration_summary.csv},
PancreSpeed 1.0 winning 96 of 96 played rounds throughout the competition at mean composite
93.735 at
\texttt{2030-pdac-1min/paper/execution/comparison/leaderboard.csv},
and a noteable T+7d Daraxonrasib postoperative restart in 29 of 32 of all the iterations at
\texttt{2030-pdac-1min/paper/execution/daraxonrasib/advisory_summary.csv}.

The exceptional processing feat is the 1001 record Phase 5 first
100 ms Arm 1 sample at
\texttt{2030-pdac-1min/paper/execution/sensors/sensor_sample_8arm.jsonl},
which no human team could review or author inline within a single
working session at the 100 kHz force sampling rate across 640
channels. The next subsection states the three persistent themes
carried forward from the present paper.

\subsection{Three Persistent Themes Carried Forward}
\label{sec:conclusions-themes}

Three themes carry forward from this PDAC paper into the next
sibling release in this repository. Clinical relevance is the
first theme: the 60 second 8 arm Whipple plus Daraxonrasib pairing
addresses the third leading cause of cancer death in the United
States and the fourth in the European Union, with 5 year overall
survival below 13 percent \cite{Siegel2025CancerStatistics}. The
2025 Dutch nationwide cohort mean ideal outcome rate of 47 percent
\cite{DutchCohort2025Whipple} is the human baseline that the 2030
PancreSpeed 1.0 target improves on, with the cross time
normalization caveat preserved in every per round LLM judge
rationale. Technical reproducibility is the second theme: the
deterministic seed contract (root seed 20260513) at
\texttt{2030-pdac-1min/paper/codegen/config/iterations.yaml}, the
doi plus url plus note triad invariant across the 35 entry
\texttt{references.bib}, and the IEC \cite{iec-80601-2-77,
iec-62304} plus FDA \cite{fda-samd, cfr-21-50-30, rev-fda-breakthrough}
plus ICH \cite{ich-e6r3} standards anchors together guarantee that
the workflow is reproducible by any third party.

Ethical alignment is the third theme: the on premises LLM
deployment, the FDA
SaMD framework opportunism \cite{fda-samd}, the TRIPOD+AI
\cite{Collins2024TRIPODAI} plus CREMLS \cite{ElEmam2024CREMLS}
reporting floor, and the United States Number 1 patient safety
framing all converge on the same alignment. Table
\ref{tab:conclusions-themes} summarizes the three themes and
their downstream sharing plan.

\clearpage

\begin{table}[H]
\caption{3 persistent themes carried forward from this PDAC full paper.}
\label{tab:conclusions-themes}
\centering
\begin{tabular}{>{\raggedright\arraybackslash}p{3.4cm} >{\raggedright\arraybackslash}p{5.5cm} >{\raggedright\arraybackslash}p{6.3cm}}
\toprule
\textbf{Theme} & \textbf{Anchor} & \textbf{Downstream sharing plan} \\
\midrule
  Clinical relevance & PDAC, 5 year OS less than 13\%, Whipple gold standard & Cohort scale up + multi center federation \\
  Technical reproducibility & Root seed 20260513, ieeetr bibstyle, doi+url+note triad & GitHub + Zenodo dual deposition \\
  Ethical alignment & On premises LLM, FDA SaMD, TRIPOD+AI, CREMLS & FDA Q-Sub + TRIPOD+AI checklist  \\
\bottomrule
\end{tabular}
\end{table}

\subsection{Safety Implications of the On-Premises LLM Advisory Layer}
\label{sec:conclusions-safety}

The on premises LLM advisory layer located
\texttt{2030-pdac-1min/paper/codegen/src/llm/compare_agent_1min.py}
and new file at
\texttt{2030-pdac-1min/paper/codegen/src/daraxonrasib/advisory.py}
minimizes significantly notable single robot error potential through second opinion
oracle isolation from the robot vendor controller.\ The mechanism
is unchanged from the v0.4.0 GBM cross study anchor
\cite{kawchak_2026_20113157}; the safety floor in the PDAC
variant is tightened by the 5 vessel safety zone gate verdicts at
\texttt{2030-pdac-1min/paper/execution/vascular/gate_verdicts.csv}
(clear 83, no_fly 6, soft_warning 6, hard_stop 5 across a 100
tick sample), by the 3 ms E stop budget at
\texttt{2030-pdac-1min/paper/codegen/src/coordination/arm_heartbeat_10khz.cpp},
and by the FDA SaMD framing \cite{fda-samd} that the Daraxonrasib
advisory layer at
\texttt{2030-pdac-1min/paper/execution/daraxonrasib/advisories.json}
preserves verbatim in every output. The combination positions the
workflow for an FDA Q-Sub pre submission as the next regulatory
step.

\subsection{Forward Path to the Final Publication}
\label{sec:conclusions-forward}

The downstream Claude Code Opus 4.7 1M Max passes that follow
this v0.9.0 full paper inherit a complete worked example. The
sibling cancer site papers (hepatocellular carcinoma,
cholangiocarcinoma, gastric cancer, and any future addition)
start by copying the v0.8.0 PDAC draft template at
\texttt{2030-pdac-1min/paper/draft-paper/} (preserved verbatim in
this PR), swapping the disease specific instruction files, and
resolving the brackets in one pass against the new disease
context. The reference inventory at
\texttt{2030-pdac-1min/paper/full-paper/references.bib} (35
entries with the doi plus url plus note triad) carries the cross
study citations that the sibling paper extends rather than
replaces, the formatting invariants at
\texttt{2030-pdac-1min/paper/full-paper/README.md} carry the
raggedright table column type contract plus the black text plus
single dash invariants, and the per section structure (abstract,
introduction, methods, results, discussion, limitations and
future work, conclusions, back matter) carries the page level
self standing layout convention.

Cite this article as Kawchak K. 2030: 60 Second Pancreatic Ductal Adenocarcinoma Robotic Whipple Procedure \&
Daraxonrasib Simulation. Zenodo. 2026; 10.5281/zenodo.20196639.
